\documentclass[10pt,conference]{IEEEtran}

% ---------- Packages ----------
\usepackage{graphicx}
\usepackage{amsmath,amssymb}
\usepackage{booktabs}
\usepackage{tabularx}
\usepackage{multirow}
\usepackage{xcolor}
\usepackage{enumitem}
\usepackage{microtype}
\usepackage{cite}
\usepackage{hyperref}
\usepackage[capitalize,noabbrev]{cleveref}

% IEEE-style PDF metadata and link colors.
\hypersetup{
  colorlinks=true,
  citecolor=black,
  linkcolor=black,
  urlcolor=blue,
  pdftitle={What Makes a Generalist Robot Policy?},
  pdfauthor={Tristan Meynier}
}
\setlength{\parindent}{0pt}
\setlength{\parskip}{0.6em}

% ---------- Draft helpers ----------
% Replace all \TODO{...} fields before submission.
\newcommand{\TODO}[1]{\textcolor{red!70!black}{[TODO: #1]}}
\newcommand{\pio}{\ensuremath{\pi_0}}
\newcommand{\placeholderfigure}[2]{%
  \begin{figure}[t]
    \centering
    \fbox{\parbox[c][4.2cm][c]{0.92\linewidth}{\centering #1}}
    \caption{#2}
  \end{figure}}
\newcommand{\placeholderfigurewide}[2]{%
  \begin{figure*}[t]
    \centering
    \fbox{\parbox[c][4.6cm][c]{0.95\textwidth}{\centering #1}}
    \caption{#2}
  \end{figure*}}

\title{What Makes a Generalist Robot Policy?\\
{\large A Study of Models, Tasks, and Data for Low-Cost Bimanual Manipulation}}

\author{\IEEEauthorblockN{Tristan Meynier}
\IEEEauthorblockA{Independent Researcher, Switzerland\\
\texttt{meyniertristan10@gmail.com}}}

\begin{document}

\maketitle

\begin{abstract}
Generalist robot policies aim to reuse knowledge across tasks, environments,
and embodiments while retaining reliable closed-loop control. This work studies
several ingredients of that problem using a low-cost, custom-built fixed-base
bimanual robot with 12 actuated degrees of freedom (DoF): each side has five
arm-joint DoF and one independently actuated gripper. We compare Action
Chunking Transformer (ACT) and the
pretrained vision-language-action model $\pi_{0.5}$ under matched robot-data
conditions, and perform parallel within-family ablations of simulated data and
multi-task training. For each model family, we compare real-only and
real-plus-simulation training, seven independently trained single-task
specialists and one jointly trained seven-task policy, and robustness under
standardized physical perturbations. Because raw human video does not provide
robot joint actions, the egocentric-human-data ablation is performed only for
$\pi_{0.5}$ using a cross-embodiment end-effector representation derived from
RGB-D hand tracking. We evaluate all configurations through repeated closed-loop
rollouts under matched initial conditions and report complete success, partial
success, inference latency, and 95\% confidence intervals. Recovery is evaluated
separately with controlled physical perturbations. The scope of the study is deliberately narrow:
conclusions are restricted to the custom robot embodiment, the seven benchmark tasks,
and the available data and compute budget. The robot's 3D-printable parts,
MuJoCo simulation environment, and model training and inference code are
available at \url{https://github.com/tmeynier/bimanual-models-tasks-data}.
\end{abstract}

\begin{IEEEkeywords}
bimanual manipulation, robot learning, vision-language-action models,
behavior cloning, simulation, egocentric data, closed-loop evaluation
\end{IEEEkeywords}

\placeholderfigurewide{Insert a paper overview: the custom bimanual robot hardware and the four
robot camera views on the left; the matched ACT/$\pi_{0.5}$ simulation and
multi-task ablations plus the $\pi_{0.5}$ egocentric-data ablation in the
center; example nominal and perturbed closed-loop rollouts on the right.}{Overview
of the custom bimanual platform and experimental study. ACT and $\pi_{0.5}$ are both
evaluated with real-only and real-plus-simulation training, single-task
specialists and multi-task training, and controlled perturbations. Egocentric
human data are additionally evaluated with $\pi_{0.5}$ through a
cross-embodiment end-effector representation.}

\section{Introduction}
\label{sec:introduction}

Recent robot-learning research has increasingly shifted from policies designed
for one robot and one task toward \emph{generalist robot policies}: models that
can reuse knowledge across tasks, environments, objects, and embodiments.
This direction is inspired by the foundation-model paradigm in language and
vision: training increasingly capable models on larger and more diverse
datasets, using greater computational resources, generally improves performance
and produces representations that can be adapted to many downstream tasks.
Robot learning aims to follow the same paradigm, but presents additional challenges. A robot policy must generate continuous, temporally correlated actions under partial observability, physical contact, sensor noise, and strict latency constraints. Moreover, each prediction changes the state from which the next prediction is made, so small errors can compound into unrecoverable failures. Robot-interaction data are also far more expensive to collect than text or images.

Historically, much of robot learning has treated manipulation as a collection
of task-specific problems: a policy is trained for one robot, one task, and one
demonstration distribution, without relying on a general-purpose pretrained
robot policy. Behavior cloning learns control directly from expert
demonstrations, but can suffer from distribution shift when the learned policy
visits states not covered by those demonstrations~\cite{ross2011dagger}. Action
Chunking Transformer (ACT) improves temporal consistency by predicting sequences
of future actions and has demonstrated precise bimanual manipulation using
low-cost hardware~\cite{zhao2023act}. Generative approaches such as Diffusion
Policy instead model \emph{multimodal} distributions over action trajectories,
where multimodality means that several distinct action modes can be valid for
the same observation and task~\cite{chi2023diffusion}. For example, a robot may
reach the same object from different directions or use different grasp
strategies. Modeling these alternatives as a distribution avoids collapsing
them into a single averaged action that may not correspond to any valid
behavior. These specialized
methods remain strong baselines when the task and operating conditions are well
specified, but they represent a narrower paradigm than generalist policies:
their semantic knowledge, task coverage, and ability to transfer across
embodiments are largely determined by the task-specific robot demonstrations on
which they are trained.

More recent vision-language-action (VLA) models follow a broader pretraining
paradigm. They adapt pretrained vision-language representations, learned largely
from Internet-scale image and text data, to robot control. Models in the
\pio{} family combine a pretrained vision-language backbone with a continuous
flow-matching action expert~\cite{black2024pi0}. Throughout this paper,
``\pio{} family'' refers to this policy lineage, whereas $\pi_{0.5}$ denotes the
specific member used in our experiments. In parallel, multi-task and
cross-embodiment training increase the diversity of physical experience
available to the policy. Open X-Embodiment showed positive transfer from data
collected across many robot platforms~\cite{openx2023}, while
$\pi_{0.5}$ demonstrated that heterogeneous robot, semantic, language, and web
data can support long-horizon manipulation in previously unseen
environments~\cite{black2025pi05}. Together, Internet-scale semantic pretraining
and diverse robot experience provide a route toward policies whose knowledge is
not tied to a single task or embodiment.

A natural continuation of cross-embodiment learning is to treat humans as an
additional embodiment. Egocentric human video is particularly attractive because
it can capture large numbers of manipulation tasks, objects, and environments at
a scale that is difficult to obtain with physical robots. Kareer et
al.~\cite{kareer2025humanrobot} show that human-to-robot transfer emerges when a
VLA is pretrained on sufficiently diverse robot scenes, tasks, and embodiments,
suggesting that large-scale pretraining can produce representations that are
increasingly agnostic to the acting embodiment. EgoScale further demonstrates
predictable scaling from 20,854 hours of action-labeled egocentric human video
and transfers the resulting motor priors to dexterous robot hands with limited
robot supervision~\cite{zheng2026egoscale}. These findings make egocentric human
experience a promising scalable source of cross-embodiment supervision for robot
learning.

Beyond direct policy learning, another emerging direction is the use of
\emph{world models}. A world model is a learned predictive model of environment
dynamics: given the current state or observation and, in an action-conditioned
setting, candidate actions, it predicts how the world is likely to evolve. Such
predictions can support planning by allowing an agent to compare possible future
outcomes before executing an action. Joint-Embedding Predictive Architectures
(JEPAs) learn predictive structure in representation space rather than
reconstructing every image pixel~\cite{lecun2022path}. V-JEPA~2 scales this idea
to more than one million hours of Internet video, and its action-conditioned
variant V-JEPA~2-AC is post-trained on robot trajectories to predict future
latent states under candidate actions, enabling model-based planning on physical
robots~\cite{assran2025vjepa2}. This family illustrates a complementary path
toward general-purpose robot intelligence: learning not only a policy that maps
observations to actions, but also an internal predictive model of how actions
change the physical world.

This technical report asks how several ingredients commonly associated with
more general robot policies behave in a deliberately small, resource-constrained
setting. The study uses a custom-built robot embodiment together with
behavior-cloning demonstrations, simulated trajectories, and egocentric human
data. Across seven tasks, we ask: \emph{what is the incremental effect of
simulated experience, multi-task parameter sharing, and egocentric human
supervision on closed-loop task success and robustness in bimanual
manipulation?} We additionally perform one exploratory zero-shot evaluation on
a shape-conditioned sorting task held out from all benchmark fine-tuning data.
We treat these as benchmark-specific measurements rather than evidence of
general open-world capability.

We study this question using a low-cost, custom-built fixed-base bimanual
robot with 12 actuated degrees of freedom (DoF): ten pose-controlling arm
joints and two independently actuated grippers. The benchmark contains seven tasks,
listed in approximate order of increasing difficulty:

\begin{enumerate}[leftmargin=*,label=\textbf{\arabic*.}]

    \item \textbf{Placement --- Basic pick-and-place:}
    Pick up a cube and place it inside a target container.

    \item \textbf{Picking --- Object retrieval:}
    Pick up a cube from inside a container and place it outside the container.

    \item \textbf{Sorting --- Color-conditioned sorting:}
    Sort the blue and green cubes into their respective containers.

    \item \textbf{Articulated manipulation --- Drawer opening:}
    Interpret a drawer number, locate the corresponding drawer, grasp its
    handle, and open it.

    \item \textbf{Precision manipulation --- Cube stacking:}
    Pick up the green cube and place it stably on top of the blue cube.

    \item \textbf{Bimanual collaboration --- Mid-air handover:}
    Pick up a cube from a container with the right arm, transfer it to the left arm
    in mid-air, and place it inside a second container.

    \item \textbf{Long-horizon manipulation --- Warehouse stocking:}
    Retrieve the cube of the requested color, open the
    specified drawer, place the cube inside it, and close the drawer.

\end{enumerate}

Task durations range from about 10\,s for Tasks~1--6 to 20\,s for Task~7.

We organize the experiments so that simulation and task-diversity effects are
measured within both model families. Comparisons between ACT and $\pi_{0.5}$
remain system-level comparisons rather than architecture-only ablations because
the models differ in pretraining, scale, task conditioning, and action
parameterization.

\begin{enumerate}[leftmargin=*,label=\textbf{E\arabic*.}]
    \item \textbf{Single-task model-family comparison:} ACT + real-robot BC and
    $\pi_{0.5}$ + real-robot BC are trained on Task~1 using the same robot
    demonstrations and evaluated under matched rollout conditions.

    \item \textbf{Simulation comparisons in both model families:} for ACT and
    $\pi_{0.5}$ separately, a Task~1 real-only policy is compared with
    real-plus-simulation training under two protocols. The fixed-budget
    ablation holds the total number of optimization examples constant; the
    additive experiment preserves real-data exposure while adding simulated
    examples and optimization steps.

    \item \textbf{Task-diversity ablation in both model families:} for ACT and
    $\pi_{0.5}$ separately, seven independently trained single-task specialists
    are compared with one jointly trained seven-task policy. The same real and
    simulated datasets are used per task, and the expected number of
    optimization examples per task is matched between specialist and multi-task
    conditions.

    \item \textbf{Egocentric-data ablation for $\pi_{0.5}$:} the seven-task
    $\pi_{0.5}$ policy is trained with or without aligned egocentric human data
    under the same total optimization budget and real-robot sampling rate.
    Human batches replace a matched fraction of simulation batches.

    \item \textbf{Held-out-task generalization:} the trained multi-task ACT and
    $\pi_{0.5}$ policies are evaluated without further training on a
    shape-conditioned sorting task: place a rectangular block in the left
    container and a cube in the right container. This task is excluded from all
    benchmark training and validation sets. The $\pi_{0.5}+\mathrm{Ego}$ policy
    is reported as an additional exploratory generalization condition.

    \item \textbf{Controlled perturbation evaluation:} the trained multi-task
    ACT and $\pi_{0.5}$ policies are evaluated without further training under
    the same standardized physical interventions to measure recovery separately
    from nominal task success. The $\pi_{0.5}+\mathrm{Ego}$ policy is reported
    as an additional exploratory recovery condition.
\end{enumerate}

All configurations are evaluated through matched closed-loop rollouts. Complete
and partial task success quantify nominal performance, while recovery is measured
separately with standardized physical perturbations. Offline action error is
used to select a checkpoint for each training run and is also reported as a
training diagnostic; it is not treated as the primary performance metric because
it does not capture compounding errors or closed-loop recovery.

\paragraph{Contributions.}
\begin{itemize}[leftmargin=*]
    \item We introduce a low-cost, 3D-printed bimanual robot,
    together with its MuJoCo simulation environment, four-view visual setup, and
    a seven-task benchmark.

    \item We provide matched within-family ablations for both ACT and
    $\pi_{0.5}$, measuring the effect of simulated demonstrations and multi-task
    parameter sharing while avoiding attribution of ACT--$\pi_{0.5}$ differences
    to a single architectural factor.

    \item We study whether RGB-D egocentric human demonstrations can provide
    additional supervision to $\pi_{0.5}$ through a human-hand/robot-end-effector
    representation, without fabricating robot joint labels for human video.

    \item We evaluate recovery separately from nominal success using controlled
    physical perturbations, allowing robustness differences between the trained
    policies to be measured directly.
\end{itemize}

\section{Related Work}
\label{sec:related-work}

Research toward general-purpose robot manipulation has progressed from
task-specific imitation learning to vision-language-action policies,
cross-embodiment training, human demonstrations, and predictive world models.

\paragraph{Task-specific imitation learning.}
Behavior cloning learns robot control directly from demonstrations, but errors
can compound when the policy encounters states absent from the training
data~\cite{ross2011dagger}. ACT addresses long-horizon bimanual manipulation by
predicting action chunks and temporally ensembling overlapping
predictions~\cite{zhao2023act}. Diffusion Policy and flow-matching methods
instead model distributions over action sequences, allowing multiple valid
behaviors to be represented~\cite{chi2023diffusion,lipman2023flow}.

\paragraph{Vision-language-action policies.}
Vision-language-action (VLA) models connect pretrained visual-language
representations to continuous robot control. The original $\pi_0$ combines a
vision-language backbone with a flow-matching action
expert~\cite{black2024pi0}. Its successor, $\pi_{0.5}$, introduces heterogeneous
co-training and high-level semantic supervision to improve open-world
generalization~\cite{black2025pi05}. We use $\pi_{0.5}$, rather than the original
$\pi_0$, as the common pretrained base policy in our experiments.

\paragraph{Cross-embodiment robot learning.}
Cross-embodiment learning investigates whether experience collected with one
physical system can improve control on another. This requires handling
differences in camera viewpoints, kinematics, action spaces, and hardware.
Open X-Embodiment showed that combining trajectories from heterogeneous robots can produce positive transfer across platforms~\cite{openx2023}.

\paragraph{Egocentric human experience.}
Humans can be viewed as an additional embodiment for which large and diverse
datasets are considerably easier to collect than physical robot trajectories.
Egocentric video therefore provides a scalable source of manipulation
experience, although human-to-robot transfer must bridge differences in
morphology, viewpoint, and action representation. Kareer
et al.~\cite{kareer2025humanrobot} co-fine-tune $\pi_{0.5}$ on robot and
egocentric human demonstrations; we denote this formulation as
$\pi_{0.5}+\mathrm{Ego}$. EgoScale reports predictable improvements as the
amount of action-labeled human video increases~\cite{zheng2026egoscale}.
More recently, EgoWAM shows that predicting future DINO feature
representations from egocentric data can improve out-of-distribution object and
scene generalization by up to four times~\cite{li2026egowam}. These results
suggest that human experience may provide a scalable source of
cross-embodiment supervision, particularly when the learned representation
captures embodiment-invariant scene dynamics rather than human actions alone.

\paragraph{Predictive world models.}
World models are not evaluated in our ablations, but represent a complementary
direction for future work.

Motivated by these large-scale studies, we examine whether some of these trends are observable on a custom, low-cost robot under limited data and compute.

\section{Robot Architecture}
\label{sec:robot}

\placeholderfigure{Insert three panels showing: (a) the leader arms used for
teleoperation, (b) the physical follower robot, and (c) the MuJoCo digital
twin.}{Robot platform and digital twin. The system consists of
(a) leader arms used for teleoperation and demonstration collection,
(b) the physical fixed-base bimanual follower robot, and
(c) the corresponding MuJoCo simulation environment.
\label{fig:bimanual-platform}}

\placeholderfigure{Insert an eight-panel figure arranged in two rows. The top
row shows the four physical camera views: fixed head, fixed body, left wrist,
and right wrist. The bottom row shows the corresponding MuJoCo views.}{Real
and simulated visual observations. The top row shows the four physical camera
views from the fixed head, fixed body, left wrist, and right wrist cameras. The
bottom row shows the corresponding MuJoCo viewpoints. The simulated cameras
are initialized from hardware-matched nominal poses and slightly randomized
around these poses during training.
\label{fig:camera-views}}

The platform is a low-cost, 3D-printed fixed-base bimanual robot with
12 actuated degrees of freedom (DoF):
\begin{equation*}
\begin{aligned}
N_{\mathrm{DoF}}
&= 2 \times \bigl(5~\text{arm-joint DoF}
\\
&\qquad\quad + 1~\text{gripper DoF}\bigr)
\\
&= 10~\text{arm-joint DoF}
\\
&\qquad + 2~\text{gripper DoF}
\\
&= 12.
\end{aligned}
\end{equation*}
The complete system, including the leader arms, physical follower robot, and
MuJoCo digital twin, is shown in Fig.~\ref{fig:bimanual-platform}.

\subsection{Bimanual Manipulation Arms}

The robot uses two 3D-printed arms. Each arm is actuated by six Feetech STS3215
serial-bus servomotors, following HuggingFace's LeRobot SO-101 follower-arm
configuration~\cite{lerobotso101}. Five joints control the arm configuration
and wrist orientation, while the sixth independently actuates the gripper.
Thus, each manipulator has five kinematic DoF that determine its end-effector
pose plus one gripper DoF. The gripper actuator is included in the 12-DoF system
count but does not provide an additional Cartesian pose DoF.

The serial-bus servomotors provide commanded and measured joint positions,
allowing the controller to operate the arms using joint-position feedback.
The two arms are fixed relative to the workspace; all benchmark manipulation is
therefore performed through arm and gripper motion alone. Objects and
articulated fixtures are arranged inside the reachable workspace of the
bimanual system.

\subsection{Visual Observations}

The robot uses four RGB viewpoints, shown in the top row of
Fig.~\ref{fig:camera-views}. An Intel RealSense D415 is mounted on the fixed
head above the arms and provides a global frontal view of the workspace. A
second 30\,fps USB RGB camera is mounted in the fixed robot body below the
arms, providing a lower global view. Two additional 30\,fps USB RGB cameras
are mounted on the left and right wrists and provide close-range views of
end-effector--object interactions. During robot training and deployment, the
D415 is used in RGB-only mode; its depth stream is not provided to ACT or
$\pi_{0.5}$.

The MuJoCo environment reproduces the same four nominal camera placements and
view directions as the physical robot, as shown in the bottom row of
Fig.~\ref{fig:camera-views}. Simulated observations are rendered from matched
head, body, left-wrist, and right-wrist viewpoints. Simulated camera poses are randomly perturbed around their hardware-matched
nominal poses during training to improve robustness to viewpoint variations.
Lighting and textures may also be randomized.

Before operation, the robot undergoes a calibration procedure in which the arm
joint offsets are established for both arms.

The robot's custom morphology and multi-view observation setup make embodiment
adaptation an important part of the study. Particularly this embodiment
is not represented in the original $\pi_{0.5}$ training data.

\subsection{Visual Observations}

The robot uses four RGB viewpoints. An Intel RealSense D415 is mounted on the
fixed head above the arms and provides a global frontal view of the workspace.
A second 30\,fps USB RGB camera is mounted in the fixed robot body below the
arms, providing a lower global view. Two additional 30\,fps USB RGB cameras are
mounted on the left and right wrists and provide close-range views of
end-effector/object interactions. During robot training and deployment, the
D415 is used in RGB-only mode: its depth stream is not provided to ACT or
$\pi_{0.5}$.

The MuJoCo environment reproduces the same four nominal camera placements and
view directions as the physical robot. Simulated training observations are
rendered from these matched head, body, left-wrist, and right-wrist viewpoints.
Camera extrinsics are therefore held fixed to the corresponding real-hardware
mounts in the main experiments; appearance properties such as lighting and
textures may still be randomized.

Before operation, the robot undergoes a calibration procedure in which the arm joint
offsets are established for both arms. The fixed head/body cameras are
extrinsically calibrated to the robot frame, while the wrist-camera transforms
are defined relative to their corresponding wrist links.

The robot's custom morphology and multi-view observation setup make embodiment
adaptation an important part of the study.

\section{Task Suite}
\label{sec:tasks}

The seven benchmark tasks are summarized in
Table~\ref{tab:tasks} in approximate order of increasing difficulty.

\begin{table*}[t]
\centering
\small
\caption{Manipulation tasks, ordered by approximate difficulty. All objects and fixtures are positioned within the fixed bimanual workspace.}
\label{tab:tasks}
\begin{tabularx}{\textwidth}
{@{}p{0.30\textwidth}X@{}}
\toprule
\textbf{Task title} &
\textbf{Language instruction} \\
\midrule

T1: Basic object placement &
``Pick up the cube and place it inside the target container.'' \\

T2: Object retrieval &
``Pick up the cube from the container and place it outside the container.'' \\

T3: Color-conditioned sorting &
``Sort the blue and green cubes into their respective containers.'' \\

T4: Number-conditioned drawer opening &
``Open the drawer with the requested number.'' \\

T5: Precision cube stacking &
``Place the green cube stably on top of the blue cube.'' \\

T6: Bimanual mid-air handover &
``Pick up the cube from the first container with the right arm, transfer it to the
left arm in mid-air, and place it inside the second container.'' \\

T7: Warehouse stocking &
``Retrieve the cube of the requested color, open the requested drawer, place the cube inside, and close the drawer.'' \\

\bottomrule
\end{tabularx}
\end{table*}

\section{Study Design}
\label{sec:study-design}

\subsection{Operational definition of a generalist policy}

We emphasize that, in this report, ``generalist'' is used in a restricted benchmark sense. A
single policy is considered more general within this benchmark when it can execute a
larger portion of the seven-task suite without using a separate model for each
task. We additionally report one exploratory held-out shape-sorting test, but this does
not establish broad zero-shot generalization, transfer to new environments, or
transfer to a new robot embodiment. We report
\emph{complete success}, indicating whether the full task is completed, and
\emph{partial success}, defined as the fraction of predefined subtasks completed
during the longer-horizon task. These metrics do not directly penalize collisions
or undesirable execution paths when the task is nevertheless completed, so all
rollouts are also inspected for anomalous behavior. Recovery is not inferred
from nominal success; it is evaluated separately using the controlled
perturbation protocol described in Sec.~\ref{sec:results}. For the longer-horizon
task, partial success is computed from a predefined sequence of task-specific
subgoals.

\subsection{Tested Hypotheses}

The experiments test the following benchmark-specific hypotheses. We report
both positive and negative changes and avoid treating a null result as evidence
that an ingredient is unimportant outside this setup.

\begin{description}[leftmargin=0pt,labelindent=0pt,style=nextline]
    \item[H1: Model-family comparison]
    Under the same Task~1 robot demonstration set, ACT and $\pi_{0.5}$ may
    differ in success and inference latency. Because the systems differ in
    pretraining, scale, task conditioning, and action parameterization, this
    comparison is descriptive and is not used to isolate any one factor.

    \item[H2: Simulation data]
    Within each model family, we report both a fixed-budget comparison, in which
    simulated batches replace repeated real-data batches, and an additive
    comparison, in which real-data exposure is preserved while simulation and
    optimization steps are added. The former isolates the data-source
    substitution under fixed compute; the latter measures the practical effect
    of adding simulation when additional compute is allowed.

    \item[H3: Multi-task parameter sharing]
    Within each model family, with matched per-task datasets and optimization
    exposure, a single seven-task policy may perform differently from seven
    independently trained specialists. A positive difference is interpreted as
    evidence consistent with beneficial transfer within this benchmark; a
    negative difference indicates interference.

    \item[H4: Egocentric human data]
    For $\pi_{0.5}$, under a fixed total optimization budget and fixed
    real-robot sampling rate, replacing a portion of simulation batches with
    aligned egocentric human data may change multi-task performance. This
    condition is not mirrored for standard ACT because human videos do not
    provide the robot joint-action targets expected by that baseline.

    \item[H5: Held-out-task generalization]
    Without additional fine-tuning, the multi-task $\pi_{0.5}$ policy may
    transfer to a shape-conditioned sorting instruction excluded from all
    benchmark fine-tuning data. The $\pi_{0.5}+\mathrm{Ego}$ result is reported
    as an exploratory additional condition. Standard ACT is not evaluated
    because its categorical task conditioning cannot specify an unseen task
    without learning new task-specific parameters. Any success is interpreted
    only as transfer to this held-out task, not as general open-world capability.

    \item[H6: Perturbation recovery]
    Under matched physical perturbations, ACT and $\pi_{0.5}$ may differ in
    recovery success even when their nominal success rates are similar. This is
    interpreted as robustness to the tested interventions, not as a general
    measure of self-correction. The $\pi_{0.5}+\mathrm{Ego}$ result is reported
    as an exploratory additional condition.
\end{description}

\section{Policy Models}
\label{sec:models}

\subsection{Observations, Actions, and Control}

At time \(t\), both policy families receive synchronized RGB observations from
the four robot cameras,
\begin{equation*}
\begin{aligned}
\mathbf{o}_t = \bigl\{&
\mathbf{o}^{\mathrm{head}}_t,
\mathbf{o}^{\mathrm{body}}_t,
\\
&\mathbf{o}^{\mathrm{left\ wrist}}_t,
\\
&\mathbf{o}^{\mathrm{right\ wrist}}_t
\bigr\}.
\end{aligned}
\end{equation*}
together with the robot configuration \(\mathbf{q}_t\). The D415 depth stream
is not used as a robot-policy input. $\pi_{0.5}$ is conditioned on the
natural-language instruction \(\ell\). Single-task ACT specialists require no
task identifier; the multi-task ACT model receives a learned categorical
embedding for the current task ID (T1--T7), providing a simple
task-conditioning signal without adding a pretrained language encoder.

The 12-dimensional robot configuration is ordered as
\begin{equation*}
\begin{aligned}
\mathbf{q}_t = \bigl[&
\mathbf{q}^{L,\mathrm{arm}}_t,
g^L_t,
\\
&\mathbf{q}^{R,\mathrm{arm}}_t,
g^R_t
\bigr],
\\
\mathbf{q}^{L,\mathrm{arm}}_t,
\mathbf{q}^{R,\mathrm{arm}}_t
&\in \mathbb{R}^{5},
\\
g^L_t,g^R_t &\in [0,1].
\end{aligned}
\end{equation*}
Only the ten arm-joint components enter the forward- and inverse-kinematics
calculations; the two gripper components are independent normalized actuator
states.

Both model families predict action chunks with horizon \(H=50\). ACT predicts
delta commands for the five pose-controlling joints of each arm together with
one absolute normalized gripper target per arm:
\begin{equation}
\begin{aligned}
    \pi^{\mathrm{ACT}}_\theta\!\Biggl(
    &\bigl[
    \Delta\mathbf{q}^{L,\mathrm{arm}},g^L,
    \\
    &\quad \Delta\mathbf{q}^{R,\mathrm{arm}},g^R
    \bigr]_{t:t+49}
    \\
    &\Bigm|
    \mathbf{o}_t,\mathbf{q}_t,e_{\mathrm{task}}
    \Biggr).
\end{aligned}
\end{equation}

For $\pi_{0.5}$, robot and human motion are expressed in a shared Cartesian
task space. For arm \(a\in\{L,R\}\), forward kinematics convert the measured
robot joint configuration into an end-effector pose,
\begin{equation}
    \mathbf{T}^{a}_t =
    \operatorname{FK}\!\left(\mathbf{q}^{a,\mathrm{arm}}_t\right),
\end{equation}
and the future poses in an action chunk are represented relative to the current
pose:
\begin{equation}
\begin{aligned}
    \Delta\mathbf{T}^{a}_{t,h}
    &=
    \left(\mathbf{T}^{a}_t\right)^{-1}
    \mathbf{T}^{a}_{t+h},
    \\
    h &= 0,\ldots,H-1.
\end{aligned}
\end{equation}
The egocentric human hand trajectories are converted to an analogous relative
pose representation after compensating for head-camera motion. This provides
a common geometric supervision signal without requiring a fictitious mapping
between human arm joints and robot joint angles.

At deployment, the policy predicts what each end effector should do, while the
inverse-kinematics controller determines how the robot joints should realize
that motion. Each predicted relative pose is composed with the current
end-effector pose to obtain a desired Cartesian target. The inverse-kinematics
solver is initialized from the current joint configuration and enforces joint
limits while preferring temporally continuous solutions. Because each arm has
five joints that affect its end-effector pose, an arbitrary 6-DoF target is not
always exactly reachable. The controller therefore uses a weighted Cartesian
objective that prioritizes translation and the task-relevant orientation
components, projecting the desired pose onto the arm's reachable task space.

The gripper is represented separately from the Cartesian pose. For each arm,
the raw servo position \(s\) is normalized to an absolute gripper command
\(g\in[0,1]\):
\begin{equation}
\begin{aligned}
    g &=
    \operatorname{clip}\!\Biggl(
        \frac{s-s_{\mathrm{open}}}
             {s_{\mathrm{closed}}-s_{\mathrm{open}}},
    \\
    &\qquad 0,1\Biggr),
\end{aligned}
\end{equation}
where \(s_{\mathrm{open}}\) and \(s_{\mathrm{closed}}\) are the calibrated
servo positions for that gripper, \(g=0\) denotes fully open, and \(g=1\)
denotes fully closed. The value \(g\) is a normalized actuator state rather
than an additional end-effector pose coordinate. Thus, each
$\pi_{0.5}$ action contains relative Cartesian targets for the left and right
end effectors together with one absolute gripper command per arm. Robot
demonstrations supervise all action channels. Human batches supervise only the
aligned Cartesian channels because the egocentric recordings do not provide
reliable robot-gripper targets; their gripper losses are masked rather than
assigned artificial values.

This action-space difference is one reason that direct ACT--$\pi_{0.5}$
comparisons are treated as system-level rather than architecture-only
comparisons.

At deployment, the controller operates at \(50\,\mathrm{Hz}\), and the first
\(K=25\) actions from each predicted chunk are executed before the policy
replans. This corresponds to \(0.5\,\mathrm{s}\) of control between successive
policy predictions, following the inference protocol used by the original
$\pi_0$ policy~\cite{black2024pi0}. We do not use temporal ensembling in the
main experiments because it reduced performance in preliminary tests.

\subsection{Behavior Cloning and Action Chunking}

Behavior cloning fits a policy to demonstration pairs by supervised learning.
It is simple and stable, but its training distribution is generated by the
demonstrator and generally contains mostly correct or successful behavior.
Consequently, the policy receives limited supervision for states caused by
mistakes or unexpected perturbations. During autonomous execution, small
errors can therefore move the robot into unseen states and compound over time.
Recovery demonstrations, trajectory perturbations, interactive data
aggregation, reinforcement learning, and recovery-focused fine-tuning can all
reduce this mismatch.

ACT uses a Transformer to predict action chunks and a conditional variational
autoencoder (CVAE) latent during training~\cite{zhao2023act}. The latent
variable helps represent multimodal demonstrations by capturing variations
between multiple valid action strategies. Chunking captures short-term motion
intent, reduces the effective decision horizon, and can smooth control through
temporal ensembling. ACT is comparatively compact and can be trained as either
a specialist or, in our multi-task condition, a task-ID-conditioned shared
policy.

\subsection{\texorpdfstring{\pio{} Family}{pi0 Family}: Pretrained VLA with Flow Matching}

The \pio{} family comprises vision-language-action (VLA) models that map visual
observations, language instructions, and robot state to continuous
actions~\cite{black2024pi0,black2025pi05}. The specific policy instantiated in
all experiments below is $\pi_{0.5}$.

The complete $\pi_{0.5}$ policy consists of a PaliGemma vision-language
backbone and a robotics-specific action expert. The backbone is initialized
from large-scale image--text pretraining~\cite{beyer2024paligemma}, but the
cross-embodiment prior does not come from this Internet pretraining alone.
Before adaptation to our robot, the complete $\pi_{0.5}$ checkpoint has been
jointly trained on heterogeneous data from multiple robot platforms, tasks,
camera configurations, and action spaces. Cross-embodiment knowledge therefore
enters through the shared VLA checkpoint---including both its robot-adapted
representations and action expert---rather than through a newly initialized
robot-specific output head.

The \pio{}-family action expert generates continuous action chunks using
conditional flow matching. Starting from a noisy action sample
\(\mathbf{x}_\tau\) at flow time \(\tau\in[0,1]\), the model learns a velocity
field \(v_\theta\) that transports samples from a simple noise distribution
toward the demonstrated action distribution:
\begin{equation}
\begin{aligned}
    \frac{d\mathbf{x}_\tau}{d\tau}
    &=
    v_\theta\!\Biggl(
    \mathbf{x}_\tau,\tau
    \\
    &\qquad\Bigm|
    \mathbf{o}_t,\mathbf{q}_t,\ell
    \Biggr).
\end{aligned}
\end{equation}
Because flow matching models a distribution over action sequences, it can
represent multiple valid manipulation strategies rather than averaging them
into a single action. Its main drawbacks are greater compute and memory
requirements.

\paragraph{Parameter-efficient adaptation.}
To limit training memory while preserving the pretrained cross-embodiment
representations, we initialize every VLA condition from the same complete
$\pi_{0.5}$ base checkpoint and freeze the original weights of both the
vision-language backbone and the pretrained action expert. Trainable low-rank
adaptation (LoRA) modules are inserted into the attention layers of both
components. The proprioception projection, action-input projection, and
action-output projection remain fully trainable. This fixed parameterization is
used for every VLA condition; the action expert is never fully fine-tuned.

Restricting updates to LoRA modules and the robotics-specific projection layers
reduces the risk of catastrophic forgetting, in which adaptation to the custom
robot could overwrite visual, semantic, or manipulation representations learned
during pretraining. Parameter-efficient adaptation limits rather than
eliminates this risk; the retained capabilities are therefore treated as a
pretrained prior rather than assumed to remain perfectly unchanged.

The robot-only $\pi_{0.5}$ and
$\pi_{0.5}+\mathrm{Ego}$ conditions use exactly the same trainable parameter
set, LoRA rank, optimizer, learning-rate schedule, initialization, and
checkpoint-selection procedure. The same backbone and action-expert LoRA
modules, together with the same trainable projection layers, are shared across
robot and human batches. Consequently, the egocentric ablation changes the
training data mixture without introducing additional trainable capacity.

Both policy families predict action chunks, but they differ substantially in
scale, initialization, task conditioning, output parameterization, and
adaptation procedure. The within-ACT and within-$\pi_{0.5}$ ablations therefore
provide the cleaner evidence about simulation and task diversity; direct
ACT--$\pi_{0.5}$ differences are reported only descriptively.

\section{Data}
\label{sec:data}

\subsection{Target-Robot Behavior-Cloning Data}

Real-robot demonstrations are collected on the custom bimanual robot through leader--follower
teleoperation. The fixed head D415 RGB stream, fixed body RGB stream, and both
wrist RGB streams are recorded synchronously at \(30\,\mathrm{Hz}\). The D415
depth stream is disabled for the robot-policy dataset and is not used by ACT or
$\pi_{0.5}$ during robot training or deployment. Joint positions and gripper
states are recorded at the same time. For ACT, these recorded robot states
directly define the action targets. For $\pi_{0.5}$, forward kinematics convert
the robot trajectories to relative dual-arm Cartesian end-effector trajectories,
with gripper commands retained as robot-specific action channels. Only
demonstrations that ultimately complete the task are retained, including those
that contain successful recovery behaviors.

\subsection{Simulation Data}

A MuJoCo digital twin is constructed from the robot's CAD and kinematic model.
Simulated demonstrations are generated by defining task-specific Cartesian
waypoints and using inverse kinematics to convert them into joint trajectories.
The simulated robot contains the same four observation cameras as the real
platform: fixed head, fixed body, left wrist, and right wrist. Their nominal
poses are matched to the physical camera mounts and remain fixed in the main
experiments. Object poses, visual appearance, physical parameters, and control
noise are randomized to improve sim-to-real transfer. The randomized parameters
and their sampling ranges are reported in
Appendix~\ref{app:sim-randomization}. The same simulated trajectory can
therefore provide joint-space targets for ACT and Cartesian end-effector targets
for $\pi_{0.5}$.

\subsection{Egocentric Human Data}

The egocentric dataset contains first-person RGB-D recordings of humans
performing tasks aligned with the robot benchmark. For this data collection
only, the Intel RealSense D415 depth stream is enabled. The D415 is worn in a
head-mounted configuration and records synchronized RGB and depth. MediaPipe
Hands~\cite{zhang2020mediapipehands} detects 21 hand landmarks in the RGB image.
For each visible landmark, the aligned depth measurement and calibrated camera
intrinsics are used to back-project the image coordinate to a metric 3D point.
Robust local depth statistics and temporal filtering are used to reduce missing
or noisy depth values.

A compact human ``end-effector'' frame is then constructed from the 3D hand
geometry. The palm position is derived from the wrist and metacarpophalangeal
landmarks, while palm orientation is estimated from orthogonal axes spanning
the wrist, index-base, and little-finger-base landmarks. The resulting 6-DoF
hand poses are tracked over time and converted to relative pose trajectories.
Camera ego-motion from the head-mounted recording is compensated using the
RGB-D sequence so that the pseudo-actions primarily describe hand motion rather
than head motion. The exact filtering and camera-motion parameters are reported
with the released preprocessing code.

Human videos do not contain robot joint positions or robot joint actions. We
therefore do not fabricate joint labels. Instead, following the cross-embodiment
co-training idea of Kareer et al.~\cite{kareer2025humanrobot}, we represent
human low-level actions as relative 6-DoF hand end-effector trajectories and
represent $\pi_{0.5}$ robot actions in the analogous relative end-effector
space. The human left/right hand tracks supervise the corresponding Cartesian
action channels. Human gripper labels are not synthesized; gripper-action loss
is masked on human batches and remains supervised only by robot data. Each
human trajectory also carries the task language instruction, and short subtask
descriptions may be used where available.

During $\pi_{0.5}+\mathrm{Ego}$ fine-tuning, human and robot batches are
co-trained in the same optimization run. Robot batches retain the ordinary
robot action supervision, while human batches train the shared visual-language
representation and the aligned end-effector prediction objective. Thus, the
human data can contribute manipulation and visual experience without requiring
a fictitious mapping from a human hand to the robot's joint angles. This experiment
is specific to the VLA condition and is not applied to standard ACT.

\subsection{Dataset Composition}

Table~\ref{tab:data} summarizes each training source in episodes and hours,
since adjacent video frames are highly correlated and do not represent
independent demonstrations.

\begin{table}[t]
\centering
\small
\caption{Balanced composition of the available training datasets.}
\label{tab:data}
\begin{tabular}{@{}lrrrr@{}}
\toprule
Source & Tasks & Episodes/task & Total episodes & Hours \\
\midrule
Real-robot BC    & 7 & 60  & 420   & 10.5 \\
Simulation       & 7 & 500 & 3,500 & 58.4 \\
Egocentric human & 7 & 100 & 700   & 17.5 \\
\bottomrule
\end{tabular}
\end{table}

All retained demonstrations complete their assigned task, although some include
successful recovery behaviors. The table reports the available data pool; the
main ablations do not simply consume all available samples. Instead, batch
sampling and optimization budgets are matched according to the comparison being
performed, as described in Sec.~\ref{sec:protocol}.

\section{Training and Evaluation Protocol}
\label{sec:protocol}

We retain the default hyperparameters of the official ACT and $\pi_{0.5}$
implementations unless otherwise stated. Within each model family, the
optimizer, learning rate, batch size, image preprocessing, action horizon, and
checkpoint-selection procedure are kept fixed across the relevant comparisons.
Training data are split by complete episodes to avoid frame-level leakage.
Validation mean-squared error (MSE) is measured every $\TODO{}$ steps, used to
select the checkpoint for each run, and reported as a training diagnostic; the
primary reported metric is closed-loop task success.

The experimental matrix is summarized in Table~\ref{tab:training-matrix}.
Perturbation evaluation reuses the trained multi-task checkpoints and therefore
does not require additional policy training.

\begin{table}[t]
\centering
\small
\caption{Training/evaluation matrix. Ego data are not used with standard ACT
because the human recordings lack ACT's robot joint-action targets.}
\label{tab:training-matrix}
\begin{tabular}{@{}lcc@{}}
\toprule
Condition & ACT & $\pi_{0.5}$ \\
\midrule
T1, real & Yes & Yes \\
T1, real + sim. (fixed) & Yes & Yes \\
T1, real + sim. (additive) & Yes & Yes \\
7 single-task specialists, real + simulation & Yes & Yes \\
1 seven-task model, real + simulation & Yes & Yes \\
1 seven-task model + ego & -- & Yes \\
Held-out shape sorting & Yes & Yes \\
Perturbation evaluation & Yes & Yes \\
\bottomrule
\end{tabular}
\end{table}

For the single-task ACT--$\pi_{0.5}$ comparison, both models use the same
train/validation split of Task~1 real-robot trajectories and the same evaluation
initial conditions. This remains a system-level comparison because model size,
pretraining, task conditioning, and action parameterization are not matched.

For the simulation study, we report two comparisons independently for ACT and
$\pi_{0.5}$. In the fixed-budget ablation, the real-only and
real-plus-simulation conditions use the same total number of optimization steps.
Batches assigned to the simulation fraction in the mixed condition are replaced
by resampled real-robot batches in the real-only condition; no unique real
episode is removed. The real/simulation batch ratio is $\TODO{}$.

In the additive experiment, the real-plus-simulation condition retains the same
expected number of real-robot batches as the real-only baseline and adds
simulated batches, increasing the total number of optimization steps from
$\TODO{}$ to $\TODO{}$. The fixed-budget comparison isolates the substitution
of simulated experience for repeated real-data exposure, whereas the additive
comparison measures the practical benefit of adding simulation and compute
together. Neither comparison tests simulation against collecting an equal amount
of new real-robot data.

For the task-diversity ablation, we train seven independent ACT specialists and
one task-ID-conditioned seven-task ACT policy, and separately train seven
independent $\pi_{0.5}$ specialists and one seven-task $\pi_{0.5}$ policy. Each
specialist sees only the real and simulated data for its task. The multi-task
models sample tasks uniformly. If one specialist receives \(U\) optimization
examples, the corresponding multi-task policy receives approximately \(7U\)
examples in total so that expected exposure per task is matched. This design
allows the specialist-to-multi-task effect to be measured separately for each
policy family.

For the egocentric-data ablation, only $\pi_{0.5}$ is considered. The two
seven-task VLA conditions use the same total number of optimization steps and
the same real-robot batch probability. In the human-data condition, a fraction
$\TODO{}$ of batches that would otherwise be sampled from simulation is replaced
by aligned egocentric human batches. Human samples supervise the relative
hand/end-effector action objective described in Sec.~\ref{sec:data}; missing
human gripper targets are masked rather than imputed.

The held-out-task and controlled-perturbation evaluations follow the same
general protocol. Both reuse the selected multi-task checkpoints without
additional training, use $\TODO{}$ matched rollout seeds per policy, preserve
the same camera observations, controller frequency, action-execution horizon,
and task-appropriate timeouts across policies, and report exact successes with
95\% confidence intervals.
ACT and $\pi_{0.5}$ form the primary descriptive comparison;
$\pi_{0.5}+\mathrm{Ego}$ is reported as an exploratory additional condition in
both evaluations.

For the held-out-task evaluation, we predeclare the instruction ``Place the
rectangular block in the left container and the cube in the right container''
before any policy training or test rollouts. The rectangular block, the
shape-based instruction, and the object-to-container assignment are excluded
from all real-robot, simulation, egocentric, and validation data used for
benchmark fine-tuning. The $\pi_{0.5}$ and $\pi_{0.5}+\mathrm{Ego}$ policies
receive this natural-language instruction. Because multi-task ACT uses
categorical task embeddings rather than language, it receives the existing T3
sorting-task ID; no new embedding is added or trained. The ACT--$\pi_{0.5}$
result is therefore a comparison between complete systems with different task
conditioning, not an architecture-only ablation. Since $\pi_{0.5}$ is
pretrained, ``held out'' refers specifically to the benchmark fine-tuning data
and does not assert that the task concept was absent from its original
pretraining corpus.

Complete success requires both objects to be inside their assigned containers
at the end of the rollout, with neither object in the wrong container. Partial
success is the fraction of correctly sorted objects, yielding scores of $0$,
$0.5$, or $1$. Initial object poses are sampled from the same predeclared
set for all three policies.

Checkpoint selection is based only on held-out real-robot episodes. All
closed-loop test rollouts use the selected checkpoint and are excluded from
model selection and hyperparameter tuning. We report the selected checkpoint
step, number of training seeds, total optimization examples, and data mixture
for every trained configuration.

\section{Experimental Evaluation and Results}
\label{sec:results}

We report direct model-family comparisons separately from within-family
ablations. Complete success is the primary nominal metric; partial success is
reported for the longer-horizon task, and inference latency is measured on the deployment
hardware. Success values are reported with 95\% confidence intervals.
Simulation and specialist-to-multi-task effects are computed separately for ACT
and $\pi_{0.5}$; differences between the two model families are interpreted
descriptively rather than as isolated architectural effects.

\begin{table*}[t]
\centering
\small
\caption{Closed-loop nominal results. ``Fixed'' matches total optimization
examples; ``additive'' preserves real-data exposure and adds simulation updates.
``7 specialists'' denotes seven independently trained single-task policies
evaluated on their corresponding tasks and averaged over the seven-task suite.}
\label{tab:main-results}
\begin{tabular}{@{}llccc@{}}
\toprule
Model & Training condition &
Complete success (\%) &
Partial success (\%) &
Latency (ms) \\
\midrule
ACT & T1, real
    & \TODO{} & \TODO{} & \TODO{} \\
$\pi_{0.5}$ & T1, real
    & \TODO{} & \TODO{} & \TODO{} \\
ACT & T1, real + simulation (fixed budget)
    & \TODO{} & \TODO{} & \TODO{} \\
$\pi_{0.5}$ & T1, real + simulation (fixed budget)
    & \TODO{} & \TODO{} & \TODO{} \\
ACT & T1, real + simulation (additive)
    & \TODO{} & \TODO{} & \TODO{} \\
$\pi_{0.5}$ & T1, real + simulation (additive)
    & \TODO{} & \TODO{} & \TODO{} \\
ACT & 7 specialists, real + simulation
    & \TODO{} & \TODO{} & \TODO{} \\
$\pi_{0.5}$ & 7 specialists, real + simulation
    & \TODO{} & \TODO{} & \TODO{} \\
ACT & 1 seven-task model, real + simulation
    & \TODO{} & \TODO{} & \TODO{} \\
$\pi_{0.5}$ & 1 seven-task model, real + simulation
    & \TODO{} & \TODO{} & \TODO{} \\
$\pi_{0.5}+\mathrm{Ego}$ & 1 seven-task model
    & \TODO{} & \TODO{} & \TODO{} \\
\bottomrule
\end{tabular}
\end{table*}

\paragraph{Main comparisons.}
The Task~1 ACT--$\pi_{0.5}$ comparison changes complete success by \TODO{}
percentage points and median inference latency by \TODO{} ms. We report this
only as a comparison between two complete systems and do not attribute the
difference to pretraining, architecture, model scale, or action representation
in isolation.

Under the fixed-budget protocol, replacing repeated real-data exposure with
simulation changes Task~1 success by \TODO{} percentage points for ACT and
\TODO{} percentage points for $\pi_{0.5}$. Under the additive protocol,
retaining real-data exposure and adding simulation changes success by \TODO{}
percentage points for ACT and \TODO{} percentage points for $\pi_{0.5}$, while
increasing total optimization examples from \TODO{} to \TODO{}. The first
comparison isolates data-source substitution under fixed compute; the second
measures the practical simulation-plus-compute effect.

For task diversity, the ACT specialist-to-multi-task difference is \TODO{}
percentage points and the $\pi_{0.5}$ difference is \TODO{} percentage points.
A positive difference is described as evidence consistent with beneficial
cross-task transfer within this benchmark, while a negative difference is
described as interference. Finally, replacing a matched fraction of simulation
batches with egocentric human data changes $\pi_{0.5}$ multi-task success by
\TODO{} percentage points. Any interpretation is restricted to the data mixture,
task suite, and compute budget used here.

\paragraph{Per-task results.}
Table~\ref{tab:per-task} reports complete success for both model families under
specialist and multi-task training, together with the $\pi_{0.5}+\mathrm{Ego}$
condition. This breakdown indicates whether average changes are shared across
tasks or concentrated on a small subset.

\begin{table*}[t]
\centering
\small
\caption{Closed-loop complete success rate (\%) by task.}
\label{tab:per-task}
\begin{tabular}{@{}lrrrrrrr@{}}
\toprule
Configuration & T1 & T2 & T3 & T4 & T5 & T6 & T7 \\
\midrule
ACT specialists + BC + simulation
& \TODO{} & \TODO{} & \TODO{} & \TODO{}
& \TODO{} & \TODO{} & \TODO{} \\

ACT multi-task + BC + simulation
& \TODO{} & \TODO{} & \TODO{} & \TODO{}
& \TODO{} & \TODO{} & \TODO{} \\

$\pi_{0.5}$ specialists + BC + simulation
& \TODO{} & \TODO{} & \TODO{} & \TODO{}
& \TODO{} & \TODO{} & \TODO{} \\

$\pi_{0.5}$ multi-task + BC + simulation
& \TODO{} & \TODO{} & \TODO{} & \TODO{}
& \TODO{} & \TODO{} & \TODO{} \\

$\pi_{0.5}+\mathrm{Ego}$ multi-task
& \TODO{} & \TODO{} & \TODO{} & \TODO{}
& \TODO{} & \TODO{} & \TODO{} \\
\bottomrule
\end{tabular}
\end{table*}

\paragraph{Held-out-task generalization.}
The held-out shape-sorting instruction is evaluated without fine-tuning under
matched initial conditions. Table~\ref{tab:held-out-task} reports complete
success for the same three multi-task policies used in the perturbation study.
ACT and $\pi_{0.5}$ form the primary descriptive comparison, while
$\pi_{0.5}+\mathrm{Ego}$ is an exploratory additional condition. The result is
treated as a measurement on one task rather than evidence of broad zero-shot
generalization.

\begin{table*}[t]
\centering
\small
\caption{Performance on the shape-conditioned sorting task held out from
benchmark fine-tuning. ACT and $\pi_{0.5}$ form the primary descriptive
comparison; $\pi_{0.5}+\mathrm{Ego}$ is exploratory.}
\label{tab:held-out-task}
\begin{tabular}{@{}lcc@{}}
\toprule
Policy & Successes/rollouts & Success (95\% CI) \\
\midrule
ACT multi-task + BC + simulation & \TODO{} / \TODO{} & \TODO{} \\
$\pi_{0.5}$ multi-task + BC + simulation & \TODO{} / \TODO{} & \TODO{} \\
$\pi_{0.5}+\mathrm{Ego}$ multi-task & \TODO{} / \TODO{} & \TODO{} \\
\bottomrule
\end{tabular}
\end{table*}

\paragraph{Controlled perturbation recovery.}
Nominal success does not by itself demonstrate recovery: a policy may succeed
more often simply because it makes fewer initial errors. We therefore perform a
separate perturbation evaluation using matched initial-condition seeds and the
trained multi-task checkpoints. Three standardized interventions are applied
where relevant: (i) \emph{object displacement}, in which the target cube is
moved laterally by $\TODO{}\,\mathrm{cm}$ after the arm begins its approach but
before the first grasp; (ii) \emph{post-grasp displacement}, in which a
successfully lifted cube is placed back on the workspace at a nearby reachable
pose, forcing the policy to reacquire it; and (iii) \emph{drawer reset}, in
which an opened drawer is moved $\TODO{}\,\mathrm{cm}$ toward the closed
position before the next subtask. Each perturbation is applied at a predefined
event rather than at a fixed wall-clock time.

The primary robustness comparison applies identical perturbations to the
seven-task ACT and $\pi_{0.5}$ real+simulation policies. We additionally report
the $\pi_{0.5}+\mathrm{Ego}$ checkpoint as an exploratory condition. Recovery
success is defined as completion of the original instruction after the
intervention without human assistance; the recovery gap is the difference from
the matched unperturbed completion rate.

\begin{table*}[t]
\centering
\small
\caption{Controlled perturbation recovery. Values are complete task-success
rates (\%) after intervention; the unperturbed column uses matched rollout
initializations. ACT and $\pi_{0.5}$ form the primary descriptive comparison;
$\pi_{0.5}+\mathrm{Ego}$ is exploratory.}
\label{tab:perturbation}
\begin{tabular}{@{}lcccc@{}}
\toprule
Policy & Unperturbed & Object displacement &
Post-grasp displacement & Drawer reset \\
\midrule
ACT multi-task + BC + simulation
& \TODO{} & \TODO{} & \TODO{} & \TODO{} \\
$\pi_{0.5}$ multi-task + BC + simulation
& \TODO{} & \TODO{} & \TODO{} & \TODO{} \\
$\pi_{0.5}+\mathrm{Ego}$ multi-task
& \TODO{} & \TODO{} & \TODO{} & \TODO{} \\
\bottomrule
\end{tabular}
\end{table*}

\paragraph{Qualitative observations.}
The main observed failure modes are missed grasps, incorrect object selection,
unstable placement, failed handovers, drawer-interaction errors, and loss of
visual contact caused by occlusion. We record representative recovery and
non-recovery cases from the controlled perturbation suite, but use the
quantitative perturbation metric rather than isolated examples to support
claims about recovery. Validation MSE is used for checkpoint selection and
reported as a training diagnostic, but it is not interpreted as a direct
measure of closed-loop task completion.

\section{Future Work}
\label{sec:future_work}

Future experiments should test additional robot embodiments, workspaces,
objects, and instructions, and should include learning curves over real-robot,
simulation, and egocentric-data scale. An initialization- and action-space-
matched control would help separate the effect of VLA pretraining from the
other differences between ACT and $\pi_{0.5}$. Additional data-source controls
---for example, replacing the human-data fraction with an equal amount of new
simulated or real-robot data---would further clarify whether any observed gain
is specific to human cross-embodiment experience or to additional
distributional diversity.

The existing MuJoCo environments also make reinforcement-learning fine-tuning a
natural extension. Starting from the imitation-learned ACT or $\pi_{0.5}$
checkpoints, simulated rollouts could expose the policy to off-demonstration
states and optimize task completion or recovery objectives directly. Such an
experiment would need to control reward design, exploration, and the
sim-to-real gap carefully, and is therefore left outside the present ablation
study.

A complementary direction is recovery-focused data collection. RaC (Recovery
and Correction) fine-tunes an imitation policy on human-in-the-loop
intervention trajectories in which an operator first moves the robot back
toward a familiar state and then demonstrates a corrective segment
~\cite{hu2025rac}. This is particularly relevant to the present platform because the
perturbation suite already identifies missed grasps, displaced objects, and
drawer-state errors that can be converted into targeted recovery examples.
Combining the present nominal demonstrations with RaC-style recovery and
correction data is therefore a concrete future experiment for improving
closed-loop robustness on Task~7 and the shorter manipulation tasks.

A hierarchical system combining slower semantic planning with a compact,
high-frequency controller is another direction, but is outside the scope of the
present experiments.

\section{Limitations}
\label{sec:limitations}

The study is intentionally small and its conclusions should be read within that
scope. First, all physical experiments use one custom robot, one workspace, the
same four camera placements, seven trained manipulation tasks, and one held-out
shape-sorting task. Performance on a single held-out instruction does not
establish broad zero-shot task generalization, open-world robustness, or
transfer to a second robot embodiment. The term ``generalist'' is therefore used only in the
restricted sense defined in
Sec.~\ref{sec:study-design}.

Second, the ACT--$\pi_{0.5}$ comparison is not an architecture ablation. The two
systems differ in pretraining, parameter count, visual-language representation,
task conditioning, and low-level action parameterization. The parallel ACT and
$\pi_{0.5}$ simulation and task-diversity ablations make it possible to compare
within-family effect sizes, but a difference between those effect sizes still
does not identify a single causal mechanism such as pretraining alone.
The held-out-task comparison retains this system-level interpretation: ACT uses
the trained T3 sorting-task ID, whereas $\pi_{0.5}$ receives the new
natural-language instruction.

Third, the simulation ablation compares simulated trajectories against repeated
exposure to a fixed real dataset; it does not show that simulation is preferable
to collecting an equal amount of new real-robot data. The specialist--multi-task
comparison tests shared training and parameter sharing as a package and cannot
fully separate representation sharing, regularization, task conditioning, or
optimization effects.

Fourth, the egocentric-data experiment is specific to $\pi_{0.5}$ and uses
pseudo-actions derived from RGB-D hand tracking rather than robot joint labels.
MediaPipe landmark error, RealSense depth noise, self-occlusion, camera
ego-motion estimation, and the approximation of a human hand by a 6-DoF
end-effector frame all introduce supervision noise. Human gripper actions are
not directly supervised. A performance difference therefore indicates an
effect of the chosen human-data co-training recipe, not proof that human
embodiment data will transfer in general. Because standard ACT is not trained
with the same human pseudo-action interface, the ego-data condition does not
provide an ACT--$\pi_{0.5}$ comparison.

Finally, Tasks~1--6 last roughly 10\,s and Task~7 roughly 20\,s. Task~7 is
longer within this benchmark but remains short compared with many household
or industrial long-horizon problems. Confidence intervals over physical
rollouts quantify rollout uncertainty, but unless several independently trained
seeds are evaluated they do not capture all training-run variance. The
controlled perturbations test a small set of predefined disturbances and should
not be interpreted as general self-correction capability. These limitations
motivate treating the results as evidence about this benchmark and experimental
regime rather than broad claims about general-purpose robot learning.

\section{Conclusion}
\label{sec:conclusion}

We presented a small-scale study of model choice, simulation, task diversity,
and egocentric human supervision on a low-cost fixed-base bimanual robot with
12 actuated DoF: ten arm-joint DoF and two gripper DoF. The experimental design
mirrors the simulation and multi-task ablations
within both ACT and $\pi_{0.5}$, allowing the effect of these training changes
to be measured separately for each policy family while keeping direct
ACT--$\pi_{0.5}$ comparisons descriptive. Egocentric RGB-D human data are
evaluated only for $\pi_{0.5}$ through a cross-embodiment end-effector
representation, and recovery is measured using explicit perturbations rather
than inferred from nominal success. The results should be interpreted only for
the robot embodiment and benchmark described here. \TODO{Insert the final
numerical findings using benchmark-specific language; report the within-ACT and
within-$\pi_{0.5}$ simulation and multi-task effect sizes, the ego-data effect,
the held-out-task result, and perturbation recovery with uncertainty, without
extrapolating beyond the tested tasks, data mixture, and compute budget.}

\section*{Reproducibility Statement}

The complete software stack for the custom bimanual platform is available at
\url{https://github.com/tmeynier/bimanual-models-tasks-data}. The repository includes the calibration,
demonstration-recording, policy-training, inference, and evaluation code,
together with the robot CAD files and MuJoCo simulation environment.

The arms use Feetech STS3215 servomotors, following the widely available
LeRobot SO-101 configuration. These motors and their controllers can be
purchased from standard online robotics suppliers. The platform consists of
the two fixed manipulation arms and their controllers. The released hardware
description also includes the fixed-head D415 mount, fixed body-camera mount,
and left/right wrist-camera mounts, together with the camera calibration used
to reproduce the corresponding MuJoCo viewpoints.

Datasets and trained checkpoints will be released through
\url{https://huggingface.co/tmeynier} under the following names:

\begin{itemize}[leftmargin=*]
    \item \textbf{Datasets:}
    \texttt{bimanual-real-bc},
    \texttt{bimanual-mujoco-sim}, and
    \texttt{bimanual-egocentric}.

    \item \textbf{Model checkpoints:}
    \texttt{bimanual-act-bc-1task},
    \texttt{bimanual-act-bc-sim-1task},
    \texttt{bimanual-act-bc-sim-task1} through
    \texttt{bimanual-act-bc-sim-task7},
    \texttt{bimanual-act-bc-sim-7task},
    \texttt{bimanual-pi05-bc-1task},
    \texttt{bimanual-pi05-bc-sim-1task},
    \texttt{bimanual-pi05-bc-sim-task1} through
    \texttt{bimanual-pi05-bc-sim-task7},
    \texttt{bimanual-pi05-bc-sim-7task}, and
    \texttt{bimanual-pi05-ego-7task}.
\end{itemize}

These releases include the training configurations, dataset metadata,
evaluation seeds, selected checkpoints, and per-rollout results required to
reproduce the experiments.

\appendix

\section{Simulation Randomization Parameters}
\label{app:sim-randomization}

Table~\ref{tab:sim-randomization} reports the parameters and sampling ranges
used to generate the simulated demonstrations.

\begin{table}[h]
\centering
\small
\caption{Domain-randomization parameters used in MuJoCo.}
\label{tab:sim-randomization}
\begin{tabular}{@{}lll@{}}
\toprule
Parameter & Distribution & Range \\
\midrule
Object position     & Uniform & \TODO{} \\
Object orientation  & Uniform & \TODO{} \\
Camera pose         & Fixed & Matched to real camera mounts \\
Lighting            & Uniform & \TODO{} \\
Object mass         & Uniform & \TODO{} \\
Contact friction    & Uniform & \TODO{} \\
Action noise        & Gaussian & \TODO{} \\
Control delay       & Uniform & \TODO{} \\
\bottomrule
\end{tabular}
\end{table}

\section{Minimum Reporting Checklist}
\begin{itemize}[leftmargin=*]
    \item Exact 12-DoF accounting: five kinematic arm joints plus one normalized
    gripper actuator per side, including the associated low-level controllers.
    \item Intrinsics/extrinsics and resolution for the fixed head, fixed body,
    left-wrist, and right-wrist cameras; confirmation that D415 depth is excluded
    from robot-policy input; observation history, action representation, chunk
    horizon, and execution horizon.
    \item Dataset episodes, hours, task/source composition, failure retention,
    train/validation/test splitting unit, and licenses/consent.
    \item Parameter counts, pretrained checkpoint, trainable layers, optimizer,
    updates, compute, seeds, and checkpoint selection.
    \item Per-task rollout counts, initial-condition seeds, randomization ranges,
    success predicates, timeouts, interventions, and safety rules.
    \item Raw per-rollout outcomes and confidence intervals, not only averages.
    \item Training seeds and matched optimization budgets for the ACT and
    $\pi_{0.5}$ simulation and task-diversity ablations.
    \item Egocentric RGB-D preprocessing details, MediaPipe confidence thresholds,
    depth filtering, camera-motion compensation, end-effector construction, and
    masked human action channels.
    \item Exact perturbation magnitudes, trigger events, applicable tasks, and
    recovery-success criteria.
\end{itemize}

\section{Suggested Supplementary Figures and Videos}
\begin{enumerate}[leftmargin=*]
    \item CAD-to-physical robot overview with the kinematic tree, all axes, and
    the four real/sim camera viewpoints.
    \item Dataset examples from real BC, simulation, and egocentric RGB-D human
    video, including MediaPipe 3D hand tracks and the derived end-effector frame.
    \item One paired rollout per ablation under the same initial-condition seed.
    \item Recovery-suite rollouts showing object displacement and grasp failure.
    \item Failure taxonomy and latency distribution on deployment hardware.
\end{enumerate}

\begin{thebibliography}{9}

\bibitem{ross2011dagger}
S.~Ross, G.~Gordon, and D.~Bagnell,
``A Reduction of Imitation Learning and Structured Prediction to No-Regret
Online Learning,'' in \emph{Proc. AISTATS}, 2011.

\bibitem{zhao2023act}
T.~Z. Zhao, V.~Kumar, S.~Levine, and C.~Finn,
``Learning Fine-Grained Bimanual Manipulation with Low-Cost Hardware,''
in \emph{Proc. Robotics: Science and Systems}, 2023.

\bibitem{chi2023diffusion}
C.~Chi, Z.~Xu, S.~Feng, E.~Cousineau, Y.~Du, B.~Burchfiel, R.~Tedrake,
and S.~Song,
``Diffusion Policy: Visuomotor Policy Learning via Action Diffusion,''
in \emph{Proc. Robotics: Science and Systems}, 2023.

\bibitem{lipman2023flow}
Y.~Lipman, R.~T. Q. Chen, H.~Ben-Hamu, M.~Nickel, and M.~Le,
``Flow Matching for Generative Modeling,''
in \emph{Proc. International Conference on Learning Representations}, 2023.

\bibitem{black2024pi0}
K.~Black et al.,
``$\pi_0$: A Vision-Language-Action Flow Model for General Robot Control,''
\emph{arXiv preprint arXiv:2410.24164}, 2024.

\bibitem{openx2023}
Open X-Embodiment Collaboration et al.,
``Open X-Embodiment: Robotic Learning Datasets and RT-X Models,''
\emph{arXiv preprint arXiv:2310.08864}, 2023.

\bibitem{black2025pi05}
K.~Black et al.,
``$\pi_{0.5}$: A Vision-Language-Action Model with Open-World
Generalization,''
\emph{arXiv preprint arXiv:2504.16054}, 2025.

\bibitem{lecun2022path}
Y.~LeCun,
``A Path Towards Autonomous Machine Intelligence,''
\emph{OpenReview}, 2022.

\bibitem{assran2025vjepa2}
M.~Assran et al.,
``V-JEPA~2: Self-Supervised Video Models Enable Understanding,
Prediction and Planning,''
\emph{arXiv preprint arXiv:2506.09985}, 2025.

\bibitem{kareer2025humanrobot}
S.~Kareer, K.~Pertsch, J.~Darpinian, J.~Hoffman, D.~Xu, S.~Levine,
C.~Finn, and S.~Nair,
``Emergence of Human to Robot Transfer in Vision-Language-Action Models,''
\emph{arXiv preprint arXiv:2512.22414}, 2025.

\bibitem{zheng2026egoscale}
R.~Zheng et al.,
``EgoScale: Scaling Dexterous Manipulation with Diverse Egocentric Human
Data,''
\emph{arXiv preprint arXiv:2602.16710}, 2026.

\bibitem{figure2026index}
Figure AI,
``Introducing Index: Building the World's Largest and Most Diverse Physical
Dataset,'' 2026. [Online]. Available:
\url{https://www.figure.ai/news/introducing-index}

\bibitem{figure2026helix25}
Figure AI,
``Helix 2.5: Zero-Shot 30-Home Generalization,'' 2026. [Online]. Available:
\url{https://www.figure.ai/news/helix-2-5-zero-shot-30-home-generalization}

\bibitem{li2026egowam}
B.~Li, X.~Yin, M.~Lin, Y.~Zhang, and D.~Xu,
``EgoWAM: World Action Models Beyond Pixels with In-the-Wild Egocentric
Human Data,''
\emph{arXiv preprint arXiv:2607.08436}, 2026. [Online]. Available:
\url{https://arxiv.org/abs/2607.08436}

\bibitem{lerobotso101}
Hugging Face,
``LeRobot SO-101 Documentation,'' 2026. [Online]. Available:
\url{https://huggingface.co/docs/lerobot/so101}


\bibitem{beyer2024paligemma}
L.~Beyer et al.,
``PaliGemma: A Versatile 3B VLM for Transfer,''
\emph{arXiv preprint arXiv:2407.07726}, 2024.

\bibitem{pavlakos2024hamer}
G.~Pavlakos, D.~Shan, I.~Radosavovic, A.~Kanazawa, D.~Fouhey, and J.~Malik,
``Reconstructing Hands in 3D with Transformers,''
in \emph{Proc. IEEE/CVF Conference on Computer Vision and Pattern Recognition},
2024.


\bibitem{zhang2020mediapipehands}
F.~Zhang, V.~Bazarevsky, A.~Vakunov, A.~Tkachenka, G.~Sung,
C.-L.~Chang, and M.~Grundmann,
``MediaPipe Hands: On-device Real-time Hand Tracking,''
\emph{arXiv preprint arXiv:2006.10214}, 2020.

\bibitem{hu2025rac}
Z.~Hu, R.~Wu, N.~Enock, J.~Li, R.~Kadakia, Z.~Erickson, and A.~Kumar,
``RaC: Robot Learning for Long-Horizon Tasks by Scaling Recovery and
Correction,''
\emph{arXiv preprint arXiv:2509.07953}, 2025.

\end{thebibliography}

\end{document}
